\documentclass[11pt,a4paper]{article}
\usepackage[utf8]{inputenc}
\usepackage[T1]{fontenc}
\usepackage[catalan]{babel}
\usepackage{lmodern}
\usepackage[sfdefault,lining]{FiraSans}
\usepackage[a4paper,top=2.3cm,bottom=2.2cm,left=1.8cm,right=1.8cm,headheight=1.3cm,headsep=0.4cm,footskip=1cm]{geometry}
\usepackage{amsmath,amssymb}
\usepackage{icomma}
\usepackage{graphicx}
\usepackage[table]{xcolor}
\usepackage{tikz}
\usepackage{fontawesome5}
\usepackage{ccicons}
\usepackage{fancyhdr}
\usepackage{lastpage}
\usepackage{tasks}
\usepackage{needspace}
\usepackage[most]{tcolorbox}
\usepackage{array}
\usepackage[hidelinks,pdftitle={Solucions de la fitxa de la unitat 2: Decimals i fraccions. 2n d'ESO},pdfsubject={Matemàtiques 2n ESO},pdfkeywords={nombres decimals, decimals exactes i periòdics, operacions amb decimals, fraccions equivalents, fracció generatriu, operacions amb fraccions, potències de fraccions, operacions combinades, problemes amb fraccions}]{hyperref}

\definecolor{navy}{HTML}{1B2748}
\definecolor{accent}{HTML}{E8603A}
\definecolor{teal}{HTML}{1F9E90}
\definecolor{soft}{HTML}{EEF1F8}
\definecolor{muted}{HTML}{5D6685}

\setlength{\parindent}{0pt}
\setlength{\parskip}{3pt}
\settasks{label=\alph*), label-format=\bfseries\color{navy}, label-width=1.3em, item-indent=1.8em, column-sep=0.6em, after-item-skip=5pt, before-skip=3pt, after-skip=4pt}

% ---------- capçalera i peu ----------
\newcommand{\llicencia}{\href{https://creativecommons.org/licenses/by-nc-sa/4.0/deed.ca}{\ccbyncsa}\ \ Llicència Creative Commons Reconeixement-NoComercial-CompartirIgual 4.0 Internacional (CC BY-NC-SA 4.0)}
\pagestyle{fancy}
\fancyhf{}
\fancyhead[L]{\raisebox{-0.35\height}{\includegraphics[height=0.95cm]{logo-gva.png}}}
\fancyhead[C]{\small\color{muted}Solucions de la fitxa · Unitat 2: Decimals i fraccions · 2n d'ESO}
\fancyhead[R]{\raisebox{-0.35\height}{\includegraphics[height=0.95cm]{logo-vermellar.png}}}
\fancyfoot[L]{\scriptsize\color{muted}\llicencia}
\fancyfoot[R]{\small\color{muted}\thepage}
\renewcommand{\headrulewidth}{0.4pt}
\renewcommand{\headrule}{\color{muted!40}\hrule height \headrulewidth}
\fancypagestyle{primera}{\fancyhf{}\renewcommand{\headrulewidth}{0pt}
  \fancyfoot[L]{\scriptsize\color{muted}\llicencia}
  \fancyfoot[R]{\small\color{muted}\thepage}}

% ---------- nivells i exercicis ----------
\newcommand{\estrelles}[1]{{\color{accent}\ifcase#1\or\faStar\faStar[regular]\faStar[regular]\or\faStar\faStar\faStar[regular]\or\faStar\faStar\faStar\fi}}
\newcounter{exer}
\newcommand{\ex}[1]{\par\needspace{4\baselineskip}\medskip\stepcounter{exer}\noindent{\color{navy}\bfseries\theexer.}\ \ {\small\estrelles{#1}}\ \ }

\newcommand{\seccio}[2]{\par\needspace{10\baselineskip}\bigskip
  \begin{tcolorbox}[colback=navy,colframe=navy,coltext=white,arc=2mm,boxrule=0pt,left=3mm,right=3mm,top=1.2mm,bottom=1.2mm]
  {\large\bfseries #1\quad #2}
  \end{tcolorbox}\vspace{-2pt}}

\newcolumntype{C}[1]{>{\centering\arraybackslash}p{#1}}
% quadre de recordatori (només a la fitxa d'exercicis)
\newcommand{\recorda}[1]{\par\smallskip\begin{tcolorbox}[colback=accent!7,colframe=accent!55,arc=2mm,boxrule=0.5pt,left=2.5mm,right=2.5mm,top=1mm,bottom=1mm]\small{\color{accent}\faLightbulb[regular]}\ \textbf{Recorda:}\ #1\end{tcolorbox}}
\newcommand{\res}[1]{\boxed{#1}}\newcommand{\rest}[1]{\fbox{#1}}

\def\fmneed{12}
% =====================================================================
% Capçalera FentMates de les fitxes: un color per a cada curs.
% Els colors són els de la pàgina inicial (LEVELS a assets/js/materials.js);
% el logo ve de _eines/logo/pdf/ (python3 _eines/logo/genera.py).
%
% Ús, al final del preàmbul de la unitat:
%   % =====================================================================
% Capçalera FentMates de les fitxes: un color per a cada curs.
% Els colors són els de la pàgina inicial (LEVELS a assets/js/materials.js);
% el logo ve de _eines/logo/pdf/ (python3 _eines/logo/genera.py).
%
% Ús, al final del preàmbul de la unitat:
%   % =====================================================================
% Capçalera FentMates de les fitxes: un color per a cada curs.
% Els colors són els de la pàgina inicial (LEVELS a assets/js/materials.js);
% el logo ve de _eines/logo/pdf/ (python3 _eines/logo/genera.py).
%
% Ús, al final del preàmbul de la unitat:
%   \input{../comu/fentmates.tex}
%   \fmcurs{1eso}
% i, al començament del document:
%   \capcalera{Matemàtiques}{1r ESO}{Fitxa d'exercicis · Unitat 2}{Divisibilitat}
% Opcional abans de \input: \def\fmneed{14} (espai mínim abans de cada secció).
% =====================================================================
\makeatletter
\def\fm@logodir{../../logo/pdf}
\providecommand{\fmneed}{10}
\def\fm@cursnom{}

% clau del curs, color i color del text damunt del color
\newcommand{\fm@curs}[3]{\@namedef{fm@color@#1}{#2}\@namedef{fm@text@#1}{#3}}
\fm@curs{1eso}{E6B84A}{navy}
\fm@curs{2eso}{3B82F6}{white}
\fm@curs{3eso}{FF7A50}{white}
\fm@curs{4eso-a}{A78BFA}{white}
\fm@curs{4eso-b}{6D28D9}{white}
\fm@curs{1batx-mi}{9B1C1C}{white}
\fm@curs{1batx-ccss}{EF6B6B}{white}
\fm@curs{2batx-mii}{1E6B47}{white}
\fm@curs{2batx-ccss}{2BB3A3}{white}

\newcommand{\fmcurs}[1]{%
  \@ifundefined{fm@color@#1}{\PackageError{fentmates}{Curs desconegut: #1}{}}{}%
  \edef\fm@tmp{\noexpand\definecolor{curs}{HTML}{\@nameuse{fm@color@#1}}%
               \noexpand\colorlet{curstext}{\@nameuse{fm@text@#1}}}\fm@tmp
  \colorlet{cursfosc}{curs!72!black}%
  \def\fm@logo{\fm@logodir/vertical-#1.pdf}%
  \def\fm@marca{\fm@logodir/marca-#1.pdf}}

% ---------- capçalera de la primera pàgina ----------
% #1 matèria (petit), #2 curs, #3 tipus de document i unitat, #4 títol
\newcommand{\capcalera}[4]{%
  \gdef\fm@cursnom{#2}%
  \par\noindent
  \begin{tcolorbox}[colback=curs,colframe=curs,boxrule=0pt,arc=2.2mm,width=3.5cm,height=2.15cm,
      valign=center,left=3mm,right=2mm,top=1mm,bottom=1mm,nobeforeafter,box align=center]
    \color{curstext}{\footnotesize #1}\par\vspace{1pt}{\fontsize{19}{21}\selectfont\bfseries #2\par}%
  \end{tcolorbox}\hspace{5mm}%
  \begin{minipage}[c]{\dimexpr\linewidth-3.5cm-5mm-2.9cm\relax}
    {\color{cursfosc}\small\bfseries\MakeUppercase{#3}}\par\vspace{2pt}
    {\color{navy}\fontsize{20}{23}\selectfont\bfseries #4\par}
  \end{minipage}\hfill
  \begin{minipage}[c]{2.6cm}\raggedleft\includegraphics[width=2.6cm]{\fm@logo}\end{minipage}\par
  \vspace{2pt}}

% ---------- barres de secció del color del curs ----------
\renewcommand{\seccio}[2]{\par\needspace{\fmneed\baselineskip}\bigskip
  \begin{tcolorbox}[colback=curs,colframe=curs,coltext=curstext,arc=2mm,boxrule=0pt,left=3mm,right=3mm,top=1.2mm,bottom=1.2mm]
  {\large\bfseries #1\quad #2}
  \end{tcolorbox}\vspace{-2pt}}

% ---------- estrelles de dificultat del color del curs (a Batxillerat no se'n posen) ----------
\@ifundefined{estrelles}{}{%
  \renewcommand{\estrelles}[1]{{\color{curs}\ifcase#1\or\faStar\faStar[regular]\faStar[regular]\or\faStar\faStar\faStar[regular]\or\faStar\faStar\faStar\fi}}}

% ---------- capçalera i peu de les pàgines ----------
% primera pàgina: logos institucionals al peu
\fancypagestyle{primera}{\fancyhf{}\renewcommand{\headrulewidth}{0pt}
  \fancyfoot[L]{\raisebox{-0.32\height}{\includegraphics[height=0.78cm]{logo-gva.png}}\hspace{5mm}%
    \raisebox{-0.32\height}{\includegraphics[height=0.78cm]{logo-vermellar.png}}}
  \fancyfoot[R]{\scriptsize\color{muted}\href{https://creativecommons.org/licenses/by-nc-sa/4.0/deed.ca}{\ccbyncsa}\ \ CC BY-NC-SA 4.0 · fentmates.org\qquad\small\thepage}}
% la resta: el símbol a l'esquerra i el curs a la dreta, amb el títol al centre com abans
\fancyhead[L]{\raisebox{-0.3\height}{\includegraphics[height=0.62cm]{\fm@marca}}}
\fancyhead[R]{\small\bfseries\color{cursfosc}\fm@cursnom}
\renewcommand{\headrule}{\color{curs!45}\hrule height \headrulewidth}
\makeatother

%   \fmcurs{1eso}
% i, al començament del document:
%   \capcalera{Matemàtiques}{1r ESO}{Fitxa d'exercicis · Unitat 2}{Divisibilitat}
% Opcional abans de \input: \def\fmneed{14} (espai mínim abans de cada secció).
% =====================================================================
\makeatletter
\def\fm@logodir{../../logo/pdf}
\providecommand{\fmneed}{10}
\def\fm@cursnom{}

% clau del curs, color i color del text damunt del color
\newcommand{\fm@curs}[3]{\@namedef{fm@color@#1}{#2}\@namedef{fm@text@#1}{#3}}
\fm@curs{1eso}{E6B84A}{navy}
\fm@curs{2eso}{3B82F6}{white}
\fm@curs{3eso}{FF7A50}{white}
\fm@curs{4eso-a}{A78BFA}{white}
\fm@curs{4eso-b}{6D28D9}{white}
\fm@curs{1batx-mi}{9B1C1C}{white}
\fm@curs{1batx-ccss}{EF6B6B}{white}
\fm@curs{2batx-mii}{1E6B47}{white}
\fm@curs{2batx-ccss}{2BB3A3}{white}

\newcommand{\fmcurs}[1]{%
  \@ifundefined{fm@color@#1}{\PackageError{fentmates}{Curs desconegut: #1}{}}{}%
  \edef\fm@tmp{\noexpand\definecolor{curs}{HTML}{\@nameuse{fm@color@#1}}%
               \noexpand\colorlet{curstext}{\@nameuse{fm@text@#1}}}\fm@tmp
  \colorlet{cursfosc}{curs!72!black}%
  \def\fm@logo{\fm@logodir/vertical-#1.pdf}%
  \def\fm@marca{\fm@logodir/marca-#1.pdf}}

% ---------- capçalera de la primera pàgina ----------
% #1 matèria (petit), #2 curs, #3 tipus de document i unitat, #4 títol
\newcommand{\capcalera}[4]{%
  \gdef\fm@cursnom{#2}%
  \par\noindent
  \begin{tcolorbox}[colback=curs,colframe=curs,boxrule=0pt,arc=2.2mm,width=3.5cm,height=2.15cm,
      valign=center,left=3mm,right=2mm,top=1mm,bottom=1mm,nobeforeafter,box align=center]
    \color{curstext}{\footnotesize #1}\par\vspace{1pt}{\fontsize{19}{21}\selectfont\bfseries #2\par}%
  \end{tcolorbox}\hspace{5mm}%
  \begin{minipage}[c]{\dimexpr\linewidth-3.5cm-5mm-2.9cm\relax}
    {\color{cursfosc}\small\bfseries\MakeUppercase{#3}}\par\vspace{2pt}
    {\color{navy}\fontsize{20}{23}\selectfont\bfseries #4\par}
  \end{minipage}\hfill
  \begin{minipage}[c]{2.6cm}\raggedleft\includegraphics[width=2.6cm]{\fm@logo}\end{minipage}\par
  \vspace{2pt}}

% ---------- barres de secció del color del curs ----------
\renewcommand{\seccio}[2]{\par\needspace{\fmneed\baselineskip}\bigskip
  \begin{tcolorbox}[colback=curs,colframe=curs,coltext=curstext,arc=2mm,boxrule=0pt,left=3mm,right=3mm,top=1.2mm,bottom=1.2mm]
  {\large\bfseries #1\quad #2}
  \end{tcolorbox}\vspace{-2pt}}

% ---------- estrelles de dificultat del color del curs (a Batxillerat no se'n posen) ----------
\@ifundefined{estrelles}{}{%
  \renewcommand{\estrelles}[1]{{\color{curs}\ifcase#1\or\faStar\faStar[regular]\faStar[regular]\or\faStar\faStar\faStar[regular]\or\faStar\faStar\faStar\fi}}}

% ---------- capçalera i peu de les pàgines ----------
% primera pàgina: logos institucionals al peu
\fancypagestyle{primera}{\fancyhf{}\renewcommand{\headrulewidth}{0pt}
  \fancyfoot[L]{\raisebox{-0.32\height}{\includegraphics[height=0.78cm]{logo-gva.png}}\hspace{5mm}%
    \raisebox{-0.32\height}{\includegraphics[height=0.78cm]{logo-vermellar.png}}}
  \fancyfoot[R]{\scriptsize\color{muted}\href{https://creativecommons.org/licenses/by-nc-sa/4.0/deed.ca}{\ccbyncsa}\ \ CC BY-NC-SA 4.0 · fentmates.org\qquad\small\thepage}}
% la resta: el símbol a l'esquerra i el curs a la dreta, amb el títol al centre com abans
\fancyhead[L]{\raisebox{-0.3\height}{\includegraphics[height=0.62cm]{\fm@marca}}}
\fancyhead[R]{\small\bfseries\color{cursfosc}\fm@cursnom}
\renewcommand{\headrule}{\color{curs!45}\hrule height \headrulewidth}
\makeatother

%   \fmcurs{1eso}
% i, al començament del document:
%   \capcalera{Matemàtiques}{1r ESO}{Fitxa d'exercicis · Unitat 2}{Divisibilitat}
% Opcional abans de \input: \def\fmneed{14} (espai mínim abans de cada secció).
% =====================================================================
\makeatletter
\def\fm@logodir{../../logo/pdf}
\providecommand{\fmneed}{10}
\def\fm@cursnom{}

% clau del curs, color i color del text damunt del color
\newcommand{\fm@curs}[3]{\@namedef{fm@color@#1}{#2}\@namedef{fm@text@#1}{#3}}
\fm@curs{1eso}{E6B84A}{navy}
\fm@curs{2eso}{3B82F6}{white}
\fm@curs{3eso}{FF7A50}{white}
\fm@curs{4eso-a}{A78BFA}{white}
\fm@curs{4eso-b}{6D28D9}{white}
\fm@curs{1batx-mi}{9B1C1C}{white}
\fm@curs{1batx-ccss}{EF6B6B}{white}
\fm@curs{2batx-mii}{1E6B47}{white}
\fm@curs{2batx-ccss}{2BB3A3}{white}

\newcommand{\fmcurs}[1]{%
  \@ifundefined{fm@color@#1}{\PackageError{fentmates}{Curs desconegut: #1}{}}{}%
  \edef\fm@tmp{\noexpand\definecolor{curs}{HTML}{\@nameuse{fm@color@#1}}%
               \noexpand\colorlet{curstext}{\@nameuse{fm@text@#1}}}\fm@tmp
  \colorlet{cursfosc}{curs!72!black}%
  \def\fm@logo{\fm@logodir/vertical-#1.pdf}%
  \def\fm@marca{\fm@logodir/marca-#1.pdf}}

% ---------- capçalera de la primera pàgina ----------
% #1 matèria (petit), #2 curs, #3 tipus de document i unitat, #4 títol
\newcommand{\capcalera}[4]{%
  \gdef\fm@cursnom{#2}%
  \par\noindent
  \begin{tcolorbox}[colback=curs,colframe=curs,boxrule=0pt,arc=2.2mm,width=3.5cm,height=2.15cm,
      valign=center,left=3mm,right=2mm,top=1mm,bottom=1mm,nobeforeafter,box align=center]
    \color{curstext}{\footnotesize #1}\par\vspace{1pt}{\fontsize{19}{21}\selectfont\bfseries #2\par}%
  \end{tcolorbox}\hspace{5mm}%
  \begin{minipage}[c]{\dimexpr\linewidth-3.5cm-5mm-2.9cm\relax}
    {\color{cursfosc}\small\bfseries\MakeUppercase{#3}}\par\vspace{2pt}
    {\color{navy}\fontsize{20}{23}\selectfont\bfseries #4\par}
  \end{minipage}\hfill
  \begin{minipage}[c]{2.6cm}\raggedleft\includegraphics[width=2.6cm]{\fm@logo}\end{minipage}\par
  \vspace{2pt}}

% ---------- barres de secció del color del curs ----------
\renewcommand{\seccio}[2]{\par\needspace{\fmneed\baselineskip}\bigskip
  \begin{tcolorbox}[colback=curs,colframe=curs,coltext=curstext,arc=2mm,boxrule=0pt,left=3mm,right=3mm,top=1.2mm,bottom=1.2mm]
  {\large\bfseries #1\quad #2}
  \end{tcolorbox}\vspace{-2pt}}

% ---------- estrelles de dificultat del color del curs (a Batxillerat no se'n posen) ----------
\@ifundefined{estrelles}{}{%
  \renewcommand{\estrelles}[1]{{\color{curs}\ifcase#1\or\faStar\faStar[regular]\faStar[regular]\or\faStar\faStar\faStar[regular]\or\faStar\faStar\faStar\fi}}}

% ---------- capçalera i peu de les pàgines ----------
% primera pàgina: logos institucionals al peu
\fancypagestyle{primera}{\fancyhf{}\renewcommand{\headrulewidth}{0pt}
  \fancyfoot[L]{\raisebox{-0.32\height}{\includegraphics[height=0.78cm]{logo-gva.png}}\hspace{5mm}%
    \raisebox{-0.32\height}{\includegraphics[height=0.78cm]{logo-vermellar.png}}}
  \fancyfoot[R]{\scriptsize\color{muted}\href{https://creativecommons.org/licenses/by-nc-sa/4.0/deed.ca}{\ccbyncsa}\ \ CC BY-NC-SA 4.0 · fentmates.org\qquad\small\thepage}}
% la resta: el símbol a l'esquerra i el curs a la dreta, amb el títol al centre com abans
\fancyhead[L]{\raisebox{-0.3\height}{\includegraphics[height=0.62cm]{\fm@marca}}}
\fancyhead[R]{\small\bfseries\color{cursfosc}\fm@cursnom}
\renewcommand{\headrule}{\color{curs!45}\hrule height \headrulewidth}
\makeatother

\fmcurs{2eso}

\begin{document}
\thispagestyle{primera}

\capcalera{Matemàtiques}{2n ESO}{Solucions de la fitxa · Unitat 2}{Decimals i fraccions}
\vspace{8pt}

\begin{tcolorbox}[colback=soft,colframe=soft,arc=2mm,left=3mm,right=3mm,top=1.5mm,bottom=1.5mm]
\small Intenta resoldre cada exercici abans de mirar-ne la solució. Si no t'ix el mateix resultat, repassa els passos o comprova'l amb les eines interactives dels apunts de la unitat. Les fraccions s'han de donar en forma irreductible.
\end{tcolorbox}

% =====================================================
\seccio{1}{Nombres decimals}

\ex{1} \begin{tasks}(2)
  \task $7{,}25$: 7 unitats i 25 centèsimes
  \task $0{,}408$: 408 mil·lèsimes
  \task $12{,}06$: 12 unitats i 6 centèsimes
  \task $3{,}5$: 3 unitats i 5 dècimes
  \task $100{,}001$: 100 unitats i 1 mil·lèsima
  \task $0{,}09$: 9 centèsimes
\end{tasks}

\ex{1} \begin{tasks}(3)
  \task 7 centèsimes $=\res{0{,}07}$
  \task 3 centèsimes $=\res{0{,}03}$
  \task 7 unitats $=\res{7}$
  \task 8 centèsimes $=\res{0{,}08}$
  \task 4 dècimes $=\res{0{,}4}$
  \task 6 mil·lèsimes $=\res{0{,}006}$
\end{tasks}

\ex{1} \begin{tasks}(3)
  \task $4{,}7\ \res{>}\ 4{,}68$
  \task $0{,}09\ \res{<}\ 0{,}1$
  \task $5{,}307\ \res{<}\ 5{,}37$
  \task $12{,}5\ \res{=}\ 12{,}5$
  \task $3{,}099\ \res{<}\ 3{,}1$
  \task $0{,}45\ \res{>}\ 0{,}405$
\end{tasks}

\ex{2} $\res{2{,}05<2{,}15<2{,}5<2{,}501<2{,}505<2{,}55}$

\ex{2} \par\smallskip{\centering\renewcommand{\arraystretch}{1.35}\begin{tabular}{|c|C{2.4cm}|C{2.4cm}|C{2.4cm}|}\hline Nombre & Truncat a les centèsimes & Arrodonit a les dècimes & Arrodonit a les centèsimes\\\hline $3{,}4682$ & $3{,}46$ & $\res{3{,}5}$ & $\res{3{,}47}$\\\hline $15{,}8962$ & $15{,}89$ & $\res{15{,}9}$ & $\res{15{,}90}$\\\hline $0{,}7049$ & $0{,}70$ & $\res{0{,}7}$ & $\res{0{,}70}$\\\hline $29{,}996$ & $29{,}99$ & $\res{30{,}0}$ & $\res{30{,}00}$\\\hline\end{tabular}\par}

\ex{3} Per exemple $4{,}701$; $4{,}705$; $4{,}7099$ (n'hi ha infinits). \ El més xicotet és $\res{7{,}50}$ i el més gran, $\res{8{,}49}$.


% =====================================================
\seccio{2}{Tipus de nombres decimals}

\ex{1} \begin{tasks}(2)
  \task $4{,}252525\ldots$: periòdic pur: $\res{4{,}\overline{25}}$
  \task $0{,}125$: \rest{exacte}
  \task $7{,}31666\ldots$: periòdic mixt: $\res{7{,}31\overline{6}}$
  \task $1{,}010010001\ldots$: \rest{no exacte i no periòdic}
  \task $0{,}777\ldots$: periòdic pur: $\res{0{,}\overline{7}}$
  \task $12{,}408408\ldots$: periòdic pur: $\res{12{,}\overline{408}}$
  \task $3{,}1252525\ldots$: periòdic mixt: $\res{3{,}1\overline{25}}$
  \task $\pi=3{,}14159\ldots$: \rest{no exacte i no periòdic}
\end{tasks}

\ex{1} \begin{tasks}(2)
  \task $5{,}\overline{6}$: part entera $\res{5}$; anteperíode no en té; període $\res{6}$
  \task $0{,}0\overline{9}$: part entera $\res{0}$; anteperíode $\res{0}$; període $\res{9}$
  \task $2{,}13\overline{4}$: part entera $\res{2}$; anteperíode $\res{13}$; període $\res{4}$
  \task $0{,}\overline{142857}$: part entera $\res{0}$; anteperíode no en té; període $\res{142857}$
\end{tasks}

\ex{2} Escrivim unes quantes xifres de cada un: $0{,}3000\ldots$; \ $0{,}3300\ldots$; \ $0{,}3333\ldots$; \ $0{,}3434\ldots$; \ $0{,}3444\ldots$\par$\res{0{,}3<0{,}33<0{,}\overline{3}<0{,}\overline{34}<0{,}3\overline{4}}$

\ex{3} Per exemple: $\res{2{,}\overline{15}}$, \ $\res{2{,}1\overline{7}}$, \ $\res{2{,}125}$ \ i \ $\res{2{,}101001000\ldots}$ (hi ha infinites respostes).


% =====================================================
\seccio{3}{Sumar, restar i multiplicar decimals}

\ex{1} \begin{tasks}(3)
  \task $23{,}45+7{,}8=\res{31{,}25}$
  \task $40-12{,}35=\res{27{,}65}$
  \task $0{,}9+0{,}35=\res{1{,}25}$
  \task $15-6{,}47=\res{8{,}53}$
  \task $3{,}06+12{,}9+0{,}4=\res{16{,}36}$
  \task $100-0{,}01=\res{99{,}99}$
\end{tasks}

\ex{1} \begin{tasks}(3)
  \task $4{,}6\cdot 3=\res{13{,}8}$
  \task $2{,}35\cdot 0{,}4=\res{0{,}94}$
  \task $0{,}6\cdot 0{,}05=\res{0{,}03}$
  \task $12{,}5\cdot 1{,}2=\res{15}$
  \task $3{,}25\cdot 1{,}4=\res{4{,}55}$
  \task $0{,}07\cdot 0{,}3=\res{0{,}021}$
\end{tasks}

\ex{1} \begin{tasks}(3)
  \task $3{,}7\cdot 100=\res{370}$
  \task $0{,}052\cdot 1000=\res{52}$
  \task $45{,}2:1000=\res{0{,}0452}$
  \task $8:100=\res{0{,}08}$
  \task $6{,}15\cdot 10=\res{61{,}5}$
  \task $0{,}3:10=\res{0{,}03}$
  \task $250\cdot 0{,}1=\res{25}$
  \task $1{,}4\cdot 0{,}01=\res{0{,}014}$
\end{tasks}

\ex{2} $2{,}5\cdot1{,}86=4{,}65$ €; \ $0{,}75\cdot14{,}20=10{,}65$ €; \ $4{,}65+10{,}65=15{,}30$ €; \ $20-15{,}30=\res{4{,}70\text{ €}}$

\ex{2} $0{,}012\cdot500=\res{6\text{ cm}}$; \ $6\cdot1000=6000\text{ cm}=\res{60\text{ m}}$

\ex{3} a) $\res{6{,}25}$ \quad b) $\res{7{,}38}$ \quad c) $\res{0{,}045}$ \quad d) $\res{0{,}3}$


% =====================================================
\seccio{4}{Dividir amb decimals}

\ex{1} \begin{tasks}(3)
  \task $8{,}4:4=\res{2{,}1}$
  \task $15:8=\res{1{,}875}$
  \task $9{,}75:3=\res{3{,}25}$
  \task $4{,}5:0{,}9=\res{5}$
  \task $2{,}73:0{,}7=\res{3{,}9}$
  \task $6:0{,}12=\res{50}$
\end{tasks}

\ex{2} \begin{tasks}(4)
  \task $10:3\approx\res{3{,}33}$
  \task $7:6\approx\res{1{,}17}$
  \task $2:7\approx\res{0{,}286}$
  \task $5{,}3:0{,}9\approx\res{5{,}89}$
\end{tasks}

\ex{2} \begin{tasks}[after-item-skip=4pt](1)
  \task $5{,}2-1{,}5\cdot 2{,}4+0{,}6:0{,}2=5{,}2-3{,}6+3=\res{4{,}6}$
  \task $\left(3{,}5-1{,}25\right)\cdot 0{,}4+2{,}7:0{,}09=2{,}25\cdot 0{,}4+30=0{,}9+30=\res{30{,}9}$
  \task $1{,}2\cdot \left(4-2{,}5\right)+0{,}8:4=1{,}2\cdot 1{,}5+0{,}2=1{,}8+0{,}2=\res{2}$
  \task $10-{\left(0{,}5+1{,}5\right)}^{2}\cdot 2=10-{2}^{2}\cdot 2=10-4\cdot 2=10-8=\res{2}$
\end{tasks}

\ex{2} $11{,}39:0{,}85=1139:85=\res{13{,}40\text{ €/kg}}$; \ $13{,}40\cdot1{,}2=\res{16{,}08\text{ €}}$

\ex{3} $31{,}5:4{,}5=\res{7\text{ L}}$ cada 100 km. \ En 260 km: $7\cdot2{,}6=18{,}2$ L; \ $18{,}2\cdot1{,}58=28{,}756\approx\res{28{,}76\text{ €}}$

\ex{3} $1{,}5:0{,}25=\res{6}$ gots i no en sobra res. \ $1{,}5:0{,}2=7{,}5$: \ $\res{7}$ gots plens i sobren $1{,}5-7\cdot0{,}2=\res{0{,}1\text{ L}}$.


% =====================================================
\seccio{5}{Fraccions}

\ex{1} \begin{tasks}(4)
  \task $\dfrac{2}{7}\cdot63=\res{18}$
  \task $\dfrac{3}{4}\cdot80=\res{60}$
  \task $\dfrac{5}{6}\cdot108=\res{90}$
  \task $\dfrac{3}{8}\cdot240=\res{90}$
\end{tasks}

\ex{1} \begin{tasks}(3)
  \task $\dfrac{84}{126}=\res{\dfrac{2}{3}}$ \ (MCD $=42$)
  \task $\dfrac{105}{180}=\res{\dfrac{7}{12}}$ \ (MCD $=15$)
  \task $\dfrac{36}{48}=\res{\dfrac{3}{4}}$ \ (MCD $=12$)
  \task $\dfrac{75}{100}=\res{\dfrac{3}{4}}$ \ (MCD $=25$)
  \task $\dfrac{240}{360}=\res{\dfrac{2}{3}}$ \ (MCD $=120$)
  \task $\dfrac{98}{140}=\res{\dfrac{7}{10}}$ \ (MCD $=14$)
\end{tasks}

\ex{2} \begin{tasks}(2)
  \task $14\cdot33=462$, \ $21\cdot22=462$: \rest{sí}
  \task $3\cdot15=45$, \ $4\cdot12=48$: \rest{no}
  \task $-6\cdot(-12)=72$, \ $8\cdot9=72$: \rest{sí}
  \task $5\cdot20=100$, \ $7\cdot15=105$: \rest{no}
\end{tasks}

\ex{2} \begin{tasks}(4)
  \task $x=\res{5}$
  \task $x=\res{-20}$
  \task $x=\res{10}$
  \task $x=\res{45}$
\end{tasks}

\ex{2} Amb denominador comú 36: $\dfrac{30}{36};\ -\dfrac{27}{36};\ \dfrac{28}{36};\ -\dfrac{24}{36};\ \dfrac{27}{36}$.\par$\res{-\dfrac{3}{4}<-\dfrac{2}{3}<\dfrac{3}{4}<\dfrac{7}{9}<\dfrac{5}{6}}$

\ex{3} Amb denominador 24: $\dfrac{8}{24}$ i $\dfrac{12}{24}$; per exemple $\res{\dfrac{9}{24},\ \dfrac{10}{24},\ \dfrac{11}{24}}$. \ Amb denominador 20: $\dfrac{8}{20}$ i $\dfrac{10}{20}$: la del mig és $\res{\dfrac{9}{20}}$.


% =====================================================
\seccio{6}{Fraccions i decimals}

\ex{1} \begin{tasks}(2)
  \task $\dfrac{9}{20}$, $20=2^{2}\cdot 5$: \rest{exacte}, $\res{0{,}45}$
  \task $\dfrac{5}{9}$, $9=3^{2}$: \rest{periòdic pur}, $\res{0{,}\overline{5}}$
  \task $\dfrac{7}{30}$, $30=2\cdot 3\cdot 5$: \rest{periòdic mixt}, $\res{0{,}2\overline{3}}$
  \task $\dfrac{21}{28}=\dfrac{3}{4}$, $4=2^{2}$: \rest{exacte}, $\res{0{,}75}$
  \task $\dfrac{4}{11}$, $11=11$: \rest{periòdic pur}, $\res{0{,}\overline{36}}$
  \task $\dfrac{7}{12}$, $12=2^{2}\cdot 3$: \rest{periòdic mixt}, $\res{0{,}58\overline{3}}$
  \task $\dfrac{13}{40}$, $40=2^{3}\cdot 5$: \rest{exacte}, $\res{0{,}325}$
  \task $\dfrac{1}{6}$, $6=2\cdot 3$: \rest{periòdic mixt}, $\res{0{,}1\overline{6}}$
\end{tasks}

\ex{1} \begin{tasks}(3)
  \task $0{,}8=\dfrac{8}{10}=\res{\dfrac{4}{5}}$
  \task $0{,}45=\dfrac{45}{100}=\res{\dfrac{9}{20}}$
  \task $1{,}2=\dfrac{12}{10}=\res{\dfrac{6}{5}}$
  \task $0{,}004=\dfrac{4}{1000}=\res{\dfrac{1}{250}}$
  \task $2{,}35=\dfrac{235}{100}=\res{\dfrac{47}{20}}$
  \task $0{,}125=\dfrac{125}{1000}=\res{\dfrac{1}{8}}$
\end{tasks}

\ex{2} \begin{tasks}[after-item-skip=4pt](2)
  \task $0{,}\overline{6}=\dfrac{6}{9}=\res{\dfrac{2}{3}}$
  \task $3{,}\overline{12}=\dfrac{312-3}{99}=\dfrac{309}{99}=\res{\dfrac{103}{33}}$
  \task $0{,}8\overline{3}=\dfrac{83-8}{90}=\dfrac{75}{90}=\res{\dfrac{5}{6}}$
  \task $2{,}1\overline{45}=\dfrac{2145-21}{990}=\dfrac{2124}{990}=\res{\dfrac{118}{55}}$
  \task $5{,}\overline{4}=\dfrac{54-5}{9}=\dfrac{49}{9}=\res{\dfrac{49}{9}}$
  \task $0{,}\overline{27}=\dfrac{27}{99}=\res{\dfrac{3}{11}}$
  \task $0{,}1\overline{6}=\dfrac{16-1}{90}=\dfrac{15}{90}=\res{\dfrac{1}{6}}$
  \task $1{,}2\overline{3}=\dfrac{123-12}{90}=\dfrac{111}{90}=\res{\dfrac{37}{30}}$
\end{tasks}

\ex{2} $37:30=1{,}2333\ldots$ \quad $100x=123{,}333\ldots$; \ $10x=12{,}333\ldots$; \ restant: $90x=111\Rightarrow x=\dfrac{111}{90}=\res{\dfrac{37}{30}}$

\ex{2} $\dfrac{2}{3}=0{,}666\ldots$, \ $\dfrac{3}{4}=0{,}75$, \ $0{,}6\overline{8}=0{,}688\ldots$: \quad $\res{\dfrac{2}{3}<0{,}67<0{,}6\overline{8}<0{,}7<\dfrac{3}{4}}$

\ex{3} a) $\dfrac{3}{9}+\dfrac{6}{9}=\res{1}$ \quad b) $\dfrac{9}{9}=\res{1}$ \quad c) $\dfrac{4}{3}\cdot\dfrac{3}{4}=\res{1}$ \quad d) $\dfrac{1}{6}:\dfrac{5}{9}=\res{\dfrac{3}{10}}$. \ En a) i b) veiem que $0{,}999\ldots$ i 1 són el mateix nombre.


% =====================================================
\seccio{7}{Suma i resta de fraccions}

\ex{1} \begin{tasks}[after-item-skip=6pt](3)
  \task $\dfrac{4}{7}+\dfrac{8}{7}=\res{\dfrac{12}{7}}$
  \task $\dfrac{3}{10}+\dfrac{2}{5}=\res{\dfrac{7}{10}}$
  \task $\dfrac{5}{6}-\dfrac{1}{4}=\res{\dfrac{7}{12}}$
  \task $\dfrac{7}{12}+\dfrac{5}{18}=\res{\dfrac{31}{36}}$
  \task $3-\dfrac{5}{4}=\res{\dfrac{7}{4}}$
  \task $1+\dfrac{2}{7}=\res{\dfrac{9}{7}}$
\end{tasks}

\ex{2} \begin{tasks}[after-item-skip=6pt](3)
  \task $\dfrac{1}{2}-\dfrac{2}{3}+\dfrac{1}{6}=\res{0}$
  \task $-\dfrac{3}{4}-\dfrac{1}{6}+2=\res{\dfrac{13}{12}}$
  \task $\dfrac{5}{6}-\dfrac{3}{4}+\dfrac{1}{3}=\res{\dfrac{5}{12}}$
  \task $\dfrac{2}{5}-\dfrac{7}{10}-1=\res{-\dfrac{13}{10}}$
  \task $-\dfrac{1}{2}-\dfrac{1}{3}-\dfrac{1}{4}=\res{-\dfrac{13}{12}}$
  \task $\dfrac{3}{8}-\dfrac{5}{12}+\dfrac{1}{6}=\res{\dfrac{1}{8}}$
\end{tasks}

\ex{2} \begin{tasks}[after-item-skip=6pt](1)
  \task $\dfrac{2}{5}-\left(\dfrac{1}{2}-\dfrac{3}{4}\right)-1=\dfrac{2}{5}-\left(-\dfrac{1}{4}\right)-1=\res{-\dfrac{7}{20}}$
  \task $\dfrac{2}{3}-\left[\dfrac{1}{4}-\left(\dfrac{5}{6}-1\right)\right]=\dfrac{2}{3}-\left[\dfrac{1}{4}-\left(-\dfrac{1}{6}\right)\right]=\dfrac{2}{3}-\dfrac{5}{12}=\res{\dfrac{1}{4}}$
  \task $\left(1-\dfrac{1}{3}\right)-\left(\dfrac{1}{2}-\dfrac{1}{5}\right)=\dfrac{2}{3}-\dfrac{3}{10}=\res{\dfrac{11}{30}}$
  \task $\dfrac{3}{4}-\left[2-\left(\dfrac{1}{3}+\dfrac{5}{6}\right)\right]=\dfrac{3}{4}-\left[2-\dfrac{7}{6}\right]=\dfrac{3}{4}-\dfrac{5}{6}=\res{-\dfrac{1}{12}}$
\end{tasks}

\ex{3} $\dfrac{5}{6}-\dfrac{3}{8}=\dfrac{20}{24}-\dfrac{9}{24}=\res{\dfrac{11}{24}}$; \qquad $\dfrac{1}{4}+\dfrac{2}{3}=\dfrac{3}{12}+\dfrac{8}{12}=\res{\dfrac{11}{12}}$

\ex{3} No es poden sumar els denominadors. Cal reduir a comú denominador: $\dfrac{10}{15}+\dfrac{3}{15}=\res{\dfrac{13}{15}}$.


% =====================================================
\seccio{8}{Multiplicació i divisió de fraccions}

\ex{1} \begin{tasks}[after-item-skip=6pt](3)
  \task $\dfrac{2}{3}\cdot \dfrac{5}{7}=\res{\dfrac{10}{21}}$
  \task $\dfrac{3}{2}\cdot \dfrac{2}{5}=\res{\dfrac{3}{5}}$
  \task $\dfrac{3}{7}\cdot 14=\res{6}$
  \task $\dfrac{4}{3}:\dfrac{5}{6}=\res{\dfrac{8}{5}}$
  \task $\dfrac{3}{5}:\dfrac{9}{10}=\res{\dfrac{2}{3}}$
  \task $6:\dfrac{3}{4}=\res{8}$
\end{tasks}

\ex{1} \begin{tasks}(5)
  \task $\res{\dfrac{5}{2}}$
  \task $\res{-\dfrac{4}{7}}$
  \task $\res{\dfrac{1}{3}}$
  \task $\res{-6}$
  \task $\res{1}$
\end{tasks}

\ex{2} \begin{tasks}[after-item-skip=6pt](3)
  \task $\left(-\dfrac{8}{15}\right)\cdot \left(-\dfrac{25}{12}\right)=\res{\dfrac{10}{9}}$
  \task $\left(-\dfrac{5}{6}\right):\dfrac{10}{9}=\res{-\dfrac{3}{4}}$
  \task $\dfrac{7}{10}:\left(-14\right)=\res{-\dfrac{1}{20}}$
  \task $\dfrac{14}{15}\cdot \dfrac{25}{21}=\res{\dfrac{10}{9}}$
  \task $\dfrac{3}{4}\cdot \dfrac{8}{9}:\dfrac{2}{3}=\res{1}$
  \task $\dfrac{2}{5}:\left(\dfrac{4}{5}\cdot \dfrac{3}{8}\right)=\res{\dfrac{4}{3}}$
\end{tasks}

\ex{2} a) $\dfrac{2}{3}\cdot\dfrac{9}{10}=\res{\dfrac{3}{5}}$ \quad b) $\dfrac{1}{2}\cdot\dfrac{4}{5}=\res{\dfrac{2}{5}}$ \quad c) $3\cdot\dfrac{5}{12}=\res{\dfrac{5}{4}}$ \quad d) $\dfrac{3}{4}:\dfrac{9}{2}=\dfrac{6}{36}=\res{\dfrac{1}{6}}$

\ex{3} $\dfrac{2}{3}:\dfrac{4}{9}=\dfrac{18}{12}=\res{\dfrac{3}{2}}$; \qquad $\dfrac{5}{8}:5=\res{\dfrac{1}{8}}$


% =====================================================
\seccio{9}{Potències de fraccions}

\ex{1} \begin{tasks}[after-item-skip=6pt](4)
  \task ${\left(\dfrac{3}{4}\right)}^{3}=\res{\dfrac{27}{64}}$
  \task ${\left(-\dfrac{2}{5}\right)}^{2}=\res{\dfrac{4}{25}}$
  \task ${\left(-\dfrac{1}{3}\right)}^{3}=\res{-\dfrac{1}{27}}$
  \task ${\left(\dfrac{7}{9}\right)}^{0}=\res{1}$
  \task ${\left(\dfrac{1}{2}\right)}^{5}=\res{\dfrac{1}{32}}$
  \task ${\left(-\dfrac{3}{2}\right)}^{4}=\res{\dfrac{81}{16}}$
  \task ${\left(-\dfrac{2}{3}\right)}^{1}=\res{-\dfrac{2}{3}}$
  \task ${\left(\dfrac{5}{2}\right)}^{2}=\res{\dfrac{25}{4}}$
\end{tasks}

\ex{1} $\res{\dfrac{4}{9}}$, \ $\res{-\dfrac{4}{9}}$, \ $\res{\dfrac{4}{3}}$, \ $\res{-\dfrac{4}{9}}$. Els parèntesis canvien el resultat!

\ex{2} \begin{tasks}[after-item-skip=8pt](2)
  \task $\left(\dfrac{2}{3}\right)^{5+2-4}=\left(\dfrac{2}{3}\right)^{3}=\res{\dfrac{8}{27}}$
  \task $\left(\dfrac{3}{5}\right)^{7-5}=\left(\dfrac{3}{5}\right)^{2}=\res{\dfrac{9}{25}}$
  \task $\left(-\dfrac{1}{2}\right)^{6-4}=\left(-\dfrac{1}{2}\right)^{2}=\res{\dfrac{1}{4}}$
  \task $\left(\dfrac{4}{7}\right)^{6-5}=\dfrac{4}{7}=\res{\dfrac{4}{7}}$
  \task $\left(\dfrac{1}{3}\right)^{2\cdot2}=\left(\dfrac{1}{3}\right)^{4}=\res{\dfrac{1}{81}}$
  \task $\left(\dfrac{5}{4}\right)^{3+0}=\left(\dfrac{5}{4}\right)^{3}=\res{\dfrac{125}{64}}$
\end{tasks}

\ex{2} \begin{tasks}[after-item-skip=6pt](2)
  \task $\left(\dfrac{5}{6}\cdot\dfrac{3}{10}\cdot4\right)^4=1^4=\res{1}$
  \task $\left(\dfrac{9}{4}:\dfrac{3}{2}\right)^3=\left(\dfrac{3}{2}\right)^3=\res{\dfrac{27}{8}}$
\end{tasks}

\ex{3} \begin{tasks}[after-item-skip=6pt](1)
  \task $\left(\dfrac{2}{3}\right)^5\cdot\left(\dfrac{2}{3}\right)^2:\left(\dfrac{2}{3}\right)^6+\left(\dfrac{1}{2}\right)^2=\left(\dfrac{2}{3}\right)^{5+2-6}+\dfrac{1}{4}=\dfrac{2}{3}+\dfrac{1}{4}=\res{\dfrac{11}{12}}$
  \task $\left[\left(\dfrac{2}{5}\right)^3\right]^2:\left(\dfrac{2}{5}\right)^4-\left(\dfrac{1}{5}\right)^2\cdot\left(\dfrac{5}{2}\right)^0=\left(\dfrac{2}{5}\right)^{6-4}-\dfrac{1}{25}\cdot1=\dfrac{4}{25}-\dfrac{1}{25}=\res{\dfrac{3}{25}}$
  \task $\left(\dfrac{6}{7}\right)^3\cdot\left(\dfrac{7}{3}\right)^3:2^2+\left(-\dfrac{1}{2}\right)^3=\left(\dfrac{6}{7}\cdot\dfrac{7}{3}\right)^3:2^2-\dfrac{1}{8}=2^3:2^2-\dfrac{1}{8}=2-\dfrac{1}{8}=\res{\dfrac{15}{8}}$
  \task $\left[\left(-\dfrac{2}{3}\right)^2\right]^3:\left(-\dfrac{2}{3}\right)^5-\left(\dfrac{1}{2}-\dfrac{3}{4}\right)^2=\left(-\dfrac{2}{3}\right)^{6-5}-\left(-\dfrac{1}{4}\right)^2=-\dfrac{2}{3}-\dfrac{1}{16}=\res{-\dfrac{35}{48}}$
  \task $\left(\dfrac{5}{4}\right)^4\cdot\left(\dfrac{8}{5}\right)^4:2^4+\left(1-\dfrac{4}{5}\right)^2\cdot5^2=\left(\dfrac{5}{4}\cdot\dfrac{8}{5}\right)^4:2^4+\left(\dfrac{1}{5}\cdot5\right)^2=2^4:2^4+1^2=1+1=\res{2}$
  \task $\left[\left(\dfrac{1}{2}\right)^3-\left(\dfrac{1}{2}\right)^2\right]:\left(-\dfrac{1}{2}\right)^3-\left(\dfrac{9}{10}\right)^3\cdot\left(\dfrac{5}{3}\right)^3:\left(\dfrac{3}{2}\right)^3=\left(\dfrac{1}{8}-\dfrac{1}{4}\right):\left(-\dfrac{1}{8}\right)-\left(\dfrac{9}{10}\cdot\dfrac{5}{3}\right)^3:\left(\dfrac{3}{2}\right)^3=\left(-\dfrac{1}{8}\right):\left(-\dfrac{1}{8}\right)-\left(\dfrac{3}{2}\right)^3:\left(\dfrac{3}{2}\right)^3=1-1=\res{0}$
\end{tasks}

\ex{3} $\left(\dfrac{2}{5}\right)^{7-6}=\res{\dfrac{2}{5}}$ \qquad $\left(-\dfrac{3}{4}\right)^{6-5}=\res{-\dfrac{3}{4}}$


% =====================================================
\seccio{10}{Operacions combinades}

\ex{1} \begin{tasks}[after-item-skip=4pt](1)
  \task $\dfrac{1}{2}+\dfrac{1}{3}\cdot \dfrac{3}{4}=\dfrac{1}{2}+\dfrac{1}{4}=\res{\dfrac{3}{4}}$
  \task $\left(\dfrac{1}{2}+\dfrac{1}{3}\right)\cdot \dfrac{3}{4}=\dfrac{5}{6}\cdot \dfrac{3}{4}=\res{\dfrac{5}{8}}$
  \task $1-\dfrac{2}{3}:\dfrac{4}{9}=1-\dfrac{3}{2}=\res{-\dfrac{1}{2}}$
  \task $\dfrac{3}{5}+\dfrac{2}{5}\cdot \dfrac{10}{3}=\dfrac{3}{5}+\dfrac{4}{3}=\res{\dfrac{29}{15}}$
  \task $\dfrac{5}{2}+\dfrac{1}{3}\cdot \dfrac{4}{5}=\dfrac{5}{2}+\dfrac{4}{15}=\res{\dfrac{83}{30}}$
  \task $\dfrac{4}{3}\cdot \left(\dfrac{2}{5}+\dfrac{1}{3}\right)=\dfrac{4}{3}\cdot \dfrac{11}{15}=\res{\dfrac{44}{45}}$
\end{tasks}

\ex{2} \begin{tasks}[after-item-skip=4pt](1)
  \task $\dfrac{1}{2}-{\left(\dfrac{2}{3}\right)}^{2}\cdot \dfrac{3}{4}=\dfrac{1}{2}-\dfrac{4}{9}\cdot \dfrac{3}{4}=\dfrac{1}{2}-\dfrac{1}{3}=\res{\dfrac{1}{6}}$
  \task ${\left(\dfrac{2}{3}\right)}^{2}-\dfrac{1}{3}\cdot \left(2-\dfrac{1}{2}\right)=\dfrac{4}{9}-\dfrac{1}{3}\cdot \dfrac{3}{2}=\dfrac{4}{9}-\dfrac{1}{2}=\res{-\dfrac{1}{18}}$
  \task $\dfrac{5}{6}:\left(\dfrac{1}{2}-\dfrac{1}{3}\right)-{2}^{2}=\dfrac{5}{6}:\dfrac{1}{6}-4=5-4=\res{1}$
  \task ${\left(-\dfrac{1}{2}\right)}^{3}+\dfrac{3}{4}:\left(1-\dfrac{1}{3}\right)=-\dfrac{1}{8}+\dfrac{3}{4}:\dfrac{2}{3}=-\dfrac{1}{8}+\dfrac{9}{8}=\res{1}$
  \task ${\left(\dfrac{3}{4}-1\right)}^{2}\cdot 8+\dfrac{1}{2}={\left(-\dfrac{1}{4}\right)}^{2}\cdot 8+\dfrac{1}{2}=\dfrac{1}{16}\cdot 8+\dfrac{1}{2}=\dfrac{1}{2}+\dfrac{1}{2}=\res{1}$
  \task $\dfrac{21}{4}-\dfrac{2}{5}:\dfrac{4}{3}=\dfrac{21}{4}-\dfrac{3}{10}=\res{\dfrac{99}{20}}$
\end{tasks}

\ex{2} \begin{tasks}[after-item-skip=4pt](1)
  \task $2-\left[\dfrac{1}{3}+{\left(\dfrac{1}{2}\right)}^{2}\right]:\dfrac{5}{6}=2-\left[\dfrac{1}{3}+\dfrac{1}{4}\right]:\dfrac{5}{6}=2-\dfrac{7}{12}:\dfrac{5}{6}=2-\dfrac{7}{10}=\res{\dfrac{13}{10}}$
  \task ${\left[\dfrac{3}{4}-{\left(\dfrac{1}{2}\right)}^{2}\right]}^{2}:{\left(-\dfrac{1}{2}\right)}^{3}+\dfrac{5}{3}\cdot \dfrac{6}{5}={\left[\dfrac{3}{4}-\dfrac{1}{4}\right]}^{2}:\left(-\dfrac{1}{8}\right)+2={\left(\dfrac{1}{2}\right)}^{2}:\left(-\dfrac{1}{8}\right)+2=\dfrac{1}{4}:\left(-\dfrac{1}{8}\right)+2=-2+2=\res{0}$
  \task $\left(\dfrac{1}{4}+\dfrac{2}{3}-\dfrac{1}{12}\right)\cdot 3+\left(\dfrac{3}{2}-\dfrac{5}{4}\right):\dfrac{1}{2}=\dfrac{5}{6}\cdot 3+\dfrac{1}{4}:\dfrac{1}{2}=\dfrac{5}{2}+\dfrac{1}{2}=\res{3}$
  \task $\dfrac{3}{2}:\dfrac{1}{4}\cdot \dfrac{1}{5}+\dfrac{3}{2}\cdot \dfrac{7}{5}=\dfrac{6}{5}+\dfrac{21}{10}=\res{\dfrac{33}{10}}$
\end{tasks}

\ex{3} $\dfrac{\frac{5}{12}}{\frac{7}{6}}-\left(-\dfrac{1}{4}\right)^2\cdot8=\dfrac{30}{84}-\dfrac{1}{16}\cdot8=\dfrac{5}{14}-\dfrac{1}{2}=\res{-\dfrac{1}{7}}$

\ex{3} a) $\res{\left(\dfrac{1}{2}+\dfrac{1}{2}\right)\cdot\dfrac{2}{3}}=\dfrac{2}{3}$ \quad b) $\res{\left(1-\dfrac{1}{3}\right):\left(\dfrac{1}{3}+1\right)}=\dfrac{2}{3}:\dfrac{4}{3}=\dfrac{1}{2}$


% =====================================================
\seccio{11}{Problemes amb fraccions}

\ex{1} $\dfrac{3}{7}\cdot28=\res{12\text{ xics}}$; \ $28-12=\res{16\text{ xiques}}$

\ex{1} $120:\dfrac{3}{5}=\dfrac{120\cdot5}{3}=\res{200\text{ pàgines}}$; \ me'n queden $\res{80}$.

\ex{1} Pa: $\res{0{,}90\text{ €}}$; \ pernil: $\res{2{,}40\text{ €}}$; \ resta: $3{,}60-0{,}90-2{,}40=\res{0{,}30\text{ €}}$

\ex{2} $\dfrac{40+112+35}{280}=\res{\dfrac{187}{280}}$; \ sobra $\res{\dfrac{93}{280}}$. \ Ha menjat més \rest{l'Elena}.

\ex{2} Dilluns: 60 L, queden 180 L. \ Dimarts: $\dfrac{2}{3}\cdot180=120$ L. \ Queden $\res{60\text{ L}}$, \ que és $\res{\dfrac{1}{4}}$ del dipòsit.

\ex{2} $\dfrac{3}{8}+\dfrac{1}{3}=\dfrac{17}{24}$; \ s'abstenen $\dfrac{7}{24}$, que són 14: \ $14:\dfrac{7}{24}=\res{48\text{ persones}}$

\ex{2} $\dfrac{3}{4}$ del líquid pesen $1{,}2-0{,}48=0{,}72$ kg; \ tot el líquid: $0{,}72:\dfrac{3}{4}=0{,}96$ kg; \ ampolla: $1{,}2-0{,}96=\res{0{,}24\text{ kg}}$

\ex{3} Cine: $\dfrac{1}{3}\cdot\dfrac{3}{5}=\dfrac{1}{5}$. \ Gasta $\dfrac{2}{5}+\dfrac{1}{5}=\dfrac{3}{5}$ i li queden $\dfrac{2}{5}$, que són 24 €: \ $24:\dfrac{2}{5}=\res{60\text{ €}}$

\ex{3} Falten $\dfrac{2}{5}\cdot\dfrac{5}{7}=\dfrac{2}{7}$ de l'etapa, que són 40 km: \ $40:\dfrac{2}{7}=\res{140\text{ km}}$. \ Al matí: $\dfrac{2}{7}\cdot140=\res{40\text{ km}}$

\ex{3} Major: $\dfrac{2}{5}\cdot180=\res{72\text{ €}}$; queden 108 €. \ Mitjà: $\dfrac{5}{8}\cdot108=\res{67{,}50\text{ €}}$. \ Menor: $108-67{,}50=\res{40{,}50\text{ €}}$, \ que és $\dfrac{40{,}5}{180}=\res{\dfrac{9}{40}}$ del total.

\ex{3} Matí: 90, queden 60. \ Vesprada: $\dfrac{3}{4}\cdot60=45$. \ Sobren $\res{15}$ croissants: $15\cdot0{,}60=\res{9\text{ €}}$


\end{document}
